\subsubsection{LightGBM Input Pipeline Comparison}
\label{sec:lightgbm-input-pipeline-comparison}

We compared two runs of LightGBM with weather features. The reference run
prepared its inputs with DuckDB. The independent run prepared Parquet files with
PySpark, then read up to 65,536 rows per batch with PyArrow into disk-backed
NumPy arrays before fitting LightGBM. Both runs used the same feature set and
chronological splits: 8,321,304 training readings from January to September
2016, 1,887,642 validation readings from October to November, and 943,555
held-out readings from December.

\begin{table}[htbp]
    \centering
    \caption{LightGBM with weather by input pipeline. Lower MAE and RMSLE indicate smaller prediction errors.}
    \label{tab:lightgbm-input-pipeline-comparison}
    \begin{tabular}{lrrrr}
        \toprule
        & \multicolumn{2}{c}{Oct.--Nov. validation}
        & \multicolumn{2}{c}{December test} \\
        \cmidrule(lr){2-3}\cmidrule(lr){4-5}
        Input pipeline & MAE & RMSLE & MAE & RMSLE \\
        \midrule
        DuckDB & \textbf{35.9330} & \textbf{0.6939} & \textbf{41.0866} & \textbf{0.6245} \\
        PySpark Parquet with batched reads & 36.3942 & 0.7017 & 41.6046 & 0.6318 \\
        \bottomrule
    \end{tabular}
\end{table}

The DuckDB-based run recorded lower errors in both evaluation periods. On the
December set, its MAE and RMSLE were 0.5180 and 0.0073 lower, respectively.
The batch-reading run successfully fitted all 200 trees, and reloading its saved
model reproduced the December scores. This shows that full-data fitting from the
Spark-prepared Parquet was feasible on the test laptop. The error differences
describe the two complete runs, with the DuckDB-based run providing the main
results for the study.
